% ============================================================
%  Emergent Gauge Structure from a 3D Brane in 4D Embedding
%  (Blank / from-scratch template with derivation placeholders)
% ============================================================
\documentclass[11pt,a4paper]{article}

% ---------- Packages ----------
\usepackage[T1]{fontenc}
\usepackage[utf8]{inputenc}
\usepackage{lmodern}
\usepackage{microtype}

\usepackage{amsmath,amssymb,amsfonts,mathtools}
\usepackage{bm}
\usepackage{physics} % bras/kets, derivatives (optional)
\usepackage{siunitx}

\usepackage{graphicx}
\usepackage{subcaption}

\usepackage{hyperref}
\hypersetup{
  colorlinks=true,
  linkcolor=blue,
  citecolor=blue,
  urlcolor=blue
}

\usepackage{geometry}
\geometry{margin=2.5cm}

% ---------- Macros ----------
\newcommand{\R}{\mathbb{R}}
\newcommand{\C}{\mathbb{C}}
\newcommand{\T}{\mathbb{T}}
\providecommand{\dd}{\mathrm{d}}
\providecommand{\ii}{\mathrm{i}}
\providecommand{\e}{\mathrm{e}}

% Indices:
%   i,j,k,... intrinsic spatial indices (1..d)
%   A,B,... embedding indices (1..N)
%   mu,nu,... spacetime indices (0..3)

\newcommand{\X}{\mathbf{X}}                 % embedding map
\newcommand{\x}{\mathbf{x}}                 % intrinsic coordinate (space)
\newcommand{\uvec}{\mathbf{u}}              % displacement field
\newcommand{\vvec}{\mathbf{v}}              % velocity field
\newcommand{\A}{\mathbf{A}}                 % gauge potential (spatial)
\newcommand{\F}{\mathbf{F}}                 % field strength / curvature
\newcommand{\Lag}{\mathcal{L}}              % Lagrangian density
\newcommand{\Ham}{\mathcal{H}}              % Hamiltonian density

% Berry / Wilczek–Zee
\newcommand{\bbA}{\mathcal{A}}              % non-Abelian connection
\newcommand{\bbF}{\mathcal{F}}              % non-Abelian curvature
\newcommand{\ab}{a}                         % U(1) Berry connection scalar components a_mu
\newcommand{\covD}{D}                       % covariant derivative symbol

% QED-like notation
\newcommand{\psibar}{\bar{\psi}}
\newcommand{\gamm}{\gamma}

% ---------- Title ----------
\title{Emergent Gauge Structure and Particle-Like Modes\\
from a 3D Brane Embedded in 4D Space}
\author{Author: \textit{(your name)}}
\date{\today}

\begin{document}
\maketitle

\begin{abstract}
This document is a from-scratch template for a brane-based deterministic wave model
in which (i) electromagnetism emerges as a U(1) gauge structure associated with a
narrowband polarization bundle, and (ii) electron-like excitations emerge as
localized multi-mode objects whose internal near-degenerate subspace carries a
non-Abelian (Wilczek--Zee) connection. The template includes (a) paper-grade
derivation steps and (b) a simulation checklist aligned with the current code
architecture (\texttt{BraneState}, \texttt{BraneGrid}, \texttt{VelocityVerletSolver},
\texttt{branesim.berry.analysis}, and the electron initialization pipeline).
\end{abstract}

\tableofcontents
\newpage

% ============================================================
\section{Scope and Claims (keep this explicit)}
\label{sec:scope}

\textbf{Goal.} Define a concrete mechanism by which an effective gauge connection
appears when the brane dynamics is projected onto a narrowband local mode subspace.
Separate what is \emph{kinematical} (connections from basis transport) from what is
\emph{dynamical} (Maxwell-like action for the connection).

\textbf{Deliverables.}
\begin{enumerate}
  \item A multiscale reduction showing the envelope field obeys covariant derivatives
        \(\covD_\mu = \partial_\mu + \ii a_\mu + \ii \bbA_\mu\).
  \item A polarization count demonstrating the EM sector yields exactly two
        gapless transverse degrees of freedom.
  \item An effective action argument that produces a Maxwell-like term
        \(f_{\mu\nu}f^{\mu\nu}\) for the U(1) curvature \(f=\dd a\).
  \item A simulation protocol that computes time- and space-like Berry connections
        from \(\uvec,\vvec\) time series and validates Maxwell identities numerically.
\end{enumerate}

% ============================================================
\section{Conceptual Model: 3D Brane Embedded in 4D Space}
\label{sec:conceptual-model}

\subsection{Kinematics and Geometry}
We model the substrate as a three-dimensional material manifold (the ``brane'')
embedded in a four-dimensional Euclidean space, with time treated as an external
parameter. Let
\[
t \in \R,\qquad
\sigma = (\sigma^1,\sigma^2,\sigma^3)\in \Omega\subset\R^3
\]
denote time and material (Lagrangian) coordinates on the brane, respectively.
The brane configuration is given by an embedding map
\[
X:\R\times\Omega \to \R^4,\qquad
(t,\sigma)\mapsto X(t,\sigma)=\big(X^1,X^2,X^3,X^4\big).
\]
Material tangent vectors are
\[
\partial_i X := \frac{\partial X}{\partial\sigma^i}\in\R^4,\qquad i=1,2,3,
\]
and they induce the metric
\[
g_{ij}(t,\sigma) := \partial_i X(t,\sigma)\cdot \partial_j X(t,\sigma),
\]
where $\cdot$ denotes the Euclidean inner product on $\R^4$.

A useful (but not required) decomposition is the static-reference form
\[
X(t,\sigma)=
\begin{pmatrix}\sigma^1\\ \sigma^2\\ \sigma^3\\ 0\end{pmatrix}
+
\begin{pmatrix}u^1(t,\sigma)\\ u^2(t,\sigma)\\ u^3(t,\sigma)\\ w(t,\sigma)\end{pmatrix},
\]
where $u=(u^1,u^2,u^3)$ is the in-brane displacement and $w=X^4$ is the
extra-dimensional displacement. The model's dimensional coupling is geometric:
all four components enter $\partial_i X$ and therefore contribute to local stretch,
elastic energy, and force.

\subsection{Directional Stretch and Geometric Coupling}
To match the structure of a nearest-neighbor spring lattice in the continuum limit,
we introduce directional stretch magnitudes
\[
a_i(t,\sigma):=\|\partial_i X(t,\sigma)\|
=\sqrt{\partial_i X\cdot\partial_i X},\qquad i=1,2,3,
\]
and a reference spacing $a_0>0$. Crucially,
\[
a_i^2=\sum_{A=1}^{4}(\partial_i X^A)^2
=\sum_{\alpha=1}^{3}(\partial_i X^\alpha)^2 + (\partial_i X^4)^2,
\]
so gradients in the extra dimension ($X^4$) modify $a_i$ and thus modify the
effective tension in \emph{all} directions. This is the fundamental coupling
mechanism between in-brane motion and extra-dimensional excitation.

\subsection{Continuous Lagrangian}
Let $\rho>0$ be the mass density with respect to material volume $d^3\sigma$.
We choose a directional stored-energy density
\[
W(\partial X) := \frac{\kappa}{2}\sum_{i=1}^{3}\big(\|\partial_i X\|-a_0\big)^2,
\]
with stiffness parameter $\kappa>0$. The action functional is
\[
S[X]=\int \! L[X]\,dt,\qquad
L[X]=\int_{\Omega}\left(\frac{\rho}{2}\|\partial_t X\|^2 - W(\partial X)\right)\,d^3\sigma,
\]
i.e.
\[
L[X]=\int_{\Omega}\left(
\frac{\rho}{2}\,\|\partial_t X\|^2
-\frac{\kappa}{2}\sum_{i=1}^{3}\big(\|\partial_i X\|-a_0\big)^2
\right)\,d^3\sigma.
\]
This Lagrangian is a direct continuum analogue of a rest-length spring network:
nonlinearity enters solely through the geometric stretch $\|\partial_i X\|$ in the
embedding space.

\subsection{Euler--Lagrange Equations}
Taking variations $X\mapsto X+\varepsilon\eta$ with compactly supported $\eta$ (or
appropriate boundary conditions), the Euler--Lagrange equations are
\[
\rho\,\partial_t^2 X
=
\sum_{i=1}^{3}\partial_i\!\left[
\kappa\left(1-\frac{a_0}{\|\partial_i X\|}\right)\partial_i X
\right].
\]
This is a nonlinear wave system in $\R^4$.
The coupling between in-brane motion and $X^4$ is explicit: the factor
$\|\partial_i X\|$ contains all four embedding components.

\subsection{Linearized Wave Limit and Calibration}
Near a homogeneous reference state with $\|\partial_i X\|\approx a_0$,
one obtains a linear wave-like regime in each embedding component,
\[
\rho\,\partial_t^2 X^A \approx T\,\Delta_{\sigma} X^A,\qquad A=1,2,3,4,
\]
with $\Delta_{\sigma}=\sum_{i=1}^3\partial_i^2$ and an effective tension scale $T$
(set by the chosen reference configuration and $\kappa$). The corresponding
characteristic wave speed is
\[
c \approx \sqrt{\frac{T}{\rho}}.
\]
In the discrete model below, the analogous calibration is expressed in terms of
$k/m$ and the lattice spacing.

\subsection{Boundary Conditions and Pre-tension}
The model supports different boundary conditions depending on the intended
experiment, e.g. periodic boundaries in one or more material directions or
Dirichlet (clamped) boundaries on $\partial\Omega$.
A pre-tensioned brane is obtained when the initial configuration is not at rest
with respect to the reference spacing (continuum: $\|\partial_i X\|\neq a_0$;
discrete: neighbor distances $\|x_m-x_n\|\neq \ell_0$). This pre-tension is then
maintained by the chosen boundary conditions.

% ============================================================
\section{Experimental Settings: Discrete Brane Lattice Model}
\label{sec:experimental-settings}

\subsection{Lattice, Degrees of Freedom, and Neighbor Topology}
For simulations we discretize the brane as a 3D lattice of nodes
\[
n=(n_1,n_2,n_3)\in \{0,\dots,N_1-1\}\times\{0,\dots,N_2-1\}\times\{0,\dots,N_3-1\}.
\]
Each node carries a position in the embedding space,
\[
x_n(t)\in\R^4,\qquad x_n=(x_n^1,x_n^2,x_n^3,x_n^4).
\]
We choose a neighbor set $\mathcal{N}(n)$, typically the $6$ axis-adjacent neighbors
(nearest neighbors) for computational efficiency. The set of undirected edges is
denoted by $E$.

\subsection{Spring Energy and Rest-length Coupling}
Each edge $(n,m)\in E$ is assigned a spring with constant $k>0$ and rest length
$\ell_0>0$. Define
\[
d_{nm}:=x_m-x_n,\qquad r_{nm}:=\|d_{nm}\|
=\sqrt{\sum_{A=1}^{4}(d_{nm}^A)^2}.
\]
The total potential energy is
\[
V(x)=\frac{k}{2}\sum_{(n,m)\in E}\big(\|x_m-x_n\|-\ell_0\big)^2.
\]
The geometric coupling between in-brane motion and the extra dimension is
automatic: writing $d_{nm}=(\Delta r,\Delta w)$ with $\Delta r\in\R^3$ and
$\Delta w\in\R$ yields
\[
r_{nm}=\sqrt{\|\Delta r\|^2+(\Delta w)^2}.
\]
Thus any excitation in $x^4$ modifies spring stretch and therefore changes forces
in the brane directions, and vice versa, without introducing additional ad-hoc
interaction terms.

\subsection{Equations of Motion}
Let $m>0$ be the point mass per node. The conservative equations of motion are
\[
m\,\ddot x_n = -\nabla_{x_n} V(x).
\]
For the spring network above, the contribution of neighbor $m\in\mathcal{N}(n)$ is
\[
F_{n\leftarrow m}
=
k\left(1-\frac{\ell_0}{\|x_m-x_n\|}\right)(x_m-x_n),
\]
and the total nodal force is the sum
\[
m\,\ddot x_n
=
\sum_{m\in\mathcal{N}(n)}
k\left(1-\frac{\ell_0}{\|x_m-x_n\|}\right)(x_m-x_n).
\]
Component-wise this reads
\[
m\,\ddot x_n^A
=
\sum_{m\in\mathcal{N}(n)}
k\left(1-\frac{\ell_0}{r_{nm}}\right)(x_m^A-x_n^A),
\qquad A=1,2,3,4.
\]
Optional numerical damping (used only for specific relaxation protocols) can be
introduced as an additional term $-\gamma\,\dot x_n$ with $\gamma\ge 0$.

\subsection{Boundary Conditions and Pre-tension in the Discrete Setting}
Periodic boundary conditions identify opposite faces of the lattice. Clamped
boundaries fix a subset of nodes to prescribed values $x_n(t)=x_n^{\mathrm{bdry}}$.
Pre-tension is implemented by choosing an initial neighbor spacing $a$ (as realized
by the initial embedding) different from the rest length $\ell_0$; the resulting
initial stress is then maintained by the boundary conditions.

\subsection{Time Integration}
We evolve the second-order system using a symplectic method such as velocity Verlet,
\[
x_n(t+\Delta t)=x_n(t)+v_n(t)\Delta t+\frac{1}{2}a_n(t)\Delta t^2,
\]
\[
a_n(t+\Delta t)=\frac{1}{m}F_n(x(t+\Delta t)),
\]
\[
v_n(t+\Delta t)=v_n(t)+\frac{1}{2}\big(a_n(t)+a_n(t+\Delta t)\big)\Delta t.
\]
The time step $\Delta t$ is chosen to resolve the fastest lattice modes implied by
$(k,m)$ and the neighbor topology.

\subsection{Discrete--Continuum Correspondence and Speed Calibration}
Let $a$ denote the lattice spacing of the initial brane configuration. A standard
order-of-magnitude mapping is
\[
m \approx \rho\,a^3,\qquad c \sim a\sqrt{\frac{k}{m}},
\]
so calibrating a target propagation speed $c$ corresponds to setting $k/m$.
Pre-tension is governed by the ratio $\ell_0/a$ and strongly affects nonlinear
regimes relevant for confinement and the emergence of effective gauge structure.

\begin{table}[t]
\centering
\caption{Brane configuration and dynamics parameters in the continuous and discrete formulations.}
\label{tab:brane-params}
\begin{tabular}{lll}
\hline
\textbf{Concept} & \textbf{Continuous model} & \textbf{Discrete lattice model} \\
\hline
Time & $t\in\R$ & $t_n = n\,\Delta t$ \\
Material domain & $\Omega\subset\R^3$ & grid indices $n=(n_1,n_2,n_3)$ \\
Material coordinate & $\sigma=(\sigma^1,\sigma^2,\sigma^3)$ &
$n\in\{0,\dots,N_1{-}1\}\times\{0,\dots,N_2{-}1\}\times\{0,\dots,N_3{-}1\}$ \\
Embedding space & $X(t,\sigma)\in\R^4$ & $x_n(t)\in\R^4$ \\
Embedding components & $X^A,\ A=1..4$ & $x_n^A,\ A=1..4$ \\
Velocity & $\partial_t X$ & $v_n=\dot x_n$ \\
Acceleration & $\partial_t^2 X$ & $a_n=\ddot x_n$ \\
\hline
Mass density / mass & $\rho$ (per $d^3\sigma$) & $m$ (per node) \\
Stiffness parameter & $\kappa$ & $k$ \\
Reference spacing / rest length & $a_0$ & $\ell_0$ \\
Initial lattice spacing (geometry) & (implicit in $X(t_0,\sigma)$) & $a$ (initial neighbor distance) \\
Directional stretch & $a_i=\|\partial_i X\|$ & $r_{nm}=\|x_m-x_n\|$ \\
Neighbor topology & directions $i=1,2,3$ & neighbor set $\mathcal{N}(n)$, edge set $E$ \\
\hline
Potential energy density &
$W=\frac{\kappa}{2}\sum_i(\|\partial_i X\|-a_0)^2$ &
$V=\frac{k}{2}\sum_{(n,m)\in E}(\|x_m-x_n\|-\ell_0)^2$ \\
Kinetic energy density & $\frac{\rho}{2}\|\partial_t X\|^2$ & $\sum_n \frac{m}{2}\|v_n\|^2$ \\
Equations of motion &
$\rho\,\partial_t^2 X=\sum_i \partial_i\!\Big[\kappa(1-\frac{a_0}{\|\partial_i X\|})\partial_i X\Big]$ &
$m\,\ddot x_n=\sum_{m\in\mathcal{N}(n)} k(1-\frac{\ell_0}{\|x_m-x_n\|})(x_m-x_n)$ \\
\hline
Wave speed calibration (typical) & $c\approx \sqrt{T/\rho}$ (effective $T$ from $\kappa,a_0$) & $c\sim a\sqrt{k/m}$ \\
Pre-tension control & $\|\partial_i X\|\neq a_0$ initially & $a\neq \ell_0$ initially \\
Boundary conditions & periodic / Dirichlet / mixed & periodic / clamped nodes / mixed \\
Time integrator & PDE integrator (e.g., symplectic) & velocity Verlet, etc. \\
\hline
\end{tabular}
\end{table}

% ============================================================
\section{Narrowband Multiscale Ansatz}
\label{sec:multiscale}

\subsection{Carrier scale and slow envelope}
Assume the excitation of interest is narrowband around a fixed angular frequency
\(\omega_0\). In the project, \(\omega_0\) is hardcoded from the Compton wavelength
\(\lambda_C\):
\[
\omega_0 = \frac{2\pi c}{\lambda_C}.
\]
Write the displacement as a slowly modulated carrier:
\begin{equation}
\uvec(\x,t)
\approx
\Re\!\left[
\e^{-\ii \omega_0 t}\sum_{a=1}^{r} \Psi^a(\x,t)\, \mathbf{e}_a(\x,t)
\right],
\label{eq:ansatz}
\end{equation}
where
\begin{itemize}
  \item \(\Psi(\x,t)\in\C^{r}\) is a complex envelope (slow in \(t\) relative to \(\omega_0^{-1}\)),
  \item \(\{\mathbf{e}_a(\x,t)\}_{a=1..r}\subset\R^4\) is a local orthonormal mode frame
        (polarization / internal basis) spanning the retained subspace,
  \item \(r=1\) corresponds to a U(1) (rank-1) bundle; \(r>1\) yields a non-Abelian bundle.
\end{itemize}

\subsection{Complex quadrature from real simulation variables}
In a real-valued simulator, construct the complex amplitude from displacement and velocity
(\texttt{branesim.berry.analysis}):
\begin{equation}
\mathbf{a}(\x,t) = \sqrt{\omega_0}\,\uvec(\x,t) + \ii\,\frac{\vvec(\x,t)}{\sqrt{\omega_0}}
\in \C^{E},
\label{eq:complex_amp}
\end{equation}
where \(E\in\{2,3,4\}\) is the number of embedding components used for the analysis.
Normalize to a ray:
\[
\ket{\psi(\x,t)} = \frac{\mathbf{a}(\x,t)}{\|\mathbf{a}(\x,t)\|}.
\]
Undefined regions are masked by amplitude thresholds.

% ============================================================
\section{Connections from Mode-Frame Transport (Berry / Wilczek--Zee)}
\label{sec:connections}

\subsection{Rank-1 (U(1)) Berry connection}
For \(r=1\), let \(\ket{u(x)}\) be the normalized local mode, with \(x^\mu=(ct,\x)\).
Define
\[
a_\mu(x) = \ii \braket{u(x)}{\partial_\mu u(x)}.
\]
The Berry phase along a path \(\Gamma\) is \(\gamma_B=\int_\Gamma a_\mu\,\dd x^\mu\),
and the curvature is
\[
f_{\mu\nu} = \partial_\mu a_\nu - \partial_\nu a_\mu.
\]

\subsection{Non-Abelian (Wilczek--Zee) connection for internal electron subspace}
For \(r>1\), choose an orthonormal basis \(\ket{u_a(x)}\), \(a=1..r\), and define
\[
\bbA_\mu^{ab}(x) = \ii \braket{u_a(x)}{\partial_\mu u_b(x)}\in\mathfrak{u}(r).
\]
Curvature:
\[
\bbF_{\mu\nu} = \partial_\mu \bbA_\nu - \partial_\nu \bbA_\mu - \ii[\bbA_\mu,\bbA_\nu].
\]

\subsection{Envelope covariant derivative}
Inserting the ansatz \eqref{eq:ansatz} and projecting onto the retained subspace produces,
generically, a covariant derivative acting on \(\Psi\):
\begin{equation}
\covD_\mu \Psi
=
\left(\partial_\mu + \ii a_\mu \mathbf{1}_r + \ii\bbA_\mu\right)\Psi.
\label{eq:covD}
\end{equation}

\paragraph{Paper-grade derivation to include (outline).}
\begin{enumerate}
  \item Write \(\uvec \approx \Re(\e^{-\ii\omega_0 t} \mathbf{E}\Psi)\) where
        \(\mathbf{E}(x)=[\mathbf{e}_1,\dots,\mathbf{e}_r]\in\R^{4\times r}\).
  \item Differentiate and project with \(\mathbf{E}^\top\) (or \(\mathbf{E}^\dagger\) in a complexified
        inner product) to obtain terms of the form \((\mathbf{E}^\top \partial_\mu \mathbf{E})\Psi\).
  \item Identify the antisymmetric part as the gauge connection; show basis changes
        \(\mathbf{E}\to \mathbf{E}U(x)\) induce the standard transformation law for \(\bbA_\mu\).
\end{enumerate}

% ============================================================
\section{QED-like Effective Lagrangian (Kinematics)}
\label{sec:qed_like}

\subsection{Local phase symmetry}
Let \(\psi(x)\) be a Dirac spinor envelope field. Under a local U(1) phase change
\(\chi(x)\):
\[
\psi \to \e^{\ii \frac{q}{\hbar}\chi(x)}\psi,
\qquad
A_\mu \to A_\mu - \partial_\mu \chi.
\]

\subsection{Minimal coupling with Berry-induced potential}
Identify the physical EM potential with the Berry connection up to units:
\[
A_\mu := \frac{\hbar}{q}\,a_\mu
\quad\Rightarrow\quad
\partial_\mu + \ii \frac{q}{\hbar}A_\mu = \partial_\mu + \ii a_\mu.
\]

\subsection{Compact (2-term) and expanded (3-term) forms}
\paragraph{2-term (compact).}
\begin{equation}
\Lag
=
\psibar\left(\ii\hbar c\,\gamm^\mu \covD_\mu - mc^2\right)\psi
-
\frac{1}{4\mu_0}F_{\mu\nu}F^{\mu\nu},
\label{eq:Lag_compact}
\end{equation}
with \(\covD_\mu=\partial_\mu + \ii\frac{q}{\hbar}A_\mu\).

\paragraph{3-term (expanded).}
\begin{equation}
\Lag
=
\underbrace{\psibar(\ii\hbar c\,\gamm^\mu \partial_\mu - mc^2)\psi}_{\text{Dirac}}
+
\underbrace{(-q\,\psibar\gamm^\mu\psi)A_\mu}_{\text{coupling}}
-
\underbrace{\frac{1}{4\mu_0}F_{\mu\nu}F^{\mu\nu}}_{\text{EM field}},
\label{eq:Lag_expanded}
\end{equation}

\subsection{Including the Wilczek--Zee sector}
If the electron envelope carries an internal index (dimension \(r\)), replace \(\psi\) with
\(\Psi\in\C^{4}\otimes\C^{r}\) (Dirac spinor valued in the internal subspace) and use
\[
\covD_\mu \Psi = \left(\partial_\mu + \ii a_\mu \mathbf{1}_r + \ii\bbA_\mu\right)\Psi.
\]
This is the place to state explicitly what \(r\) represents in the brane model
(e.g., internal polarization/standing-wave basis supporting spin-like behavior).

% ============================================================
\section{Dynamics: Where the Maxwell Term Must Come From}
\label{sec:maxwell_term}

The Berry connection provides gauge-covariant \emph{kinematics}. Electromagnetism requires
\emph{dynamics} for \(a_\mu\). In this model, the Maxwell-like term must emerge as an
effective action after coarse-graining brane degrees of freedom.

\subsection{Effective action sketch (to be made explicit)}
\begin{enumerate}
  \item Split brane variables into fast carrier-scale modes and slow envelope/mode-frame variables.
  \item Integrate out fast modes (or average over a cycle) to obtain an effective energy
        in terms of envelope gradients and frame curvature.
  \item Show the leading gauge-invariant quadratic term is \(f_{\mu\nu}f^{\mu\nu}\) with a
        computable prefactor (mapping to \(\mu_0\), \(\epsilon_0\), and an effective coupling).
\end{enumerate}

\subsection{Polarization count and suppression of extra embedding modes}
State the mechanism that restricts the EM sector to two transverse polarizations:
\begin{itemize}
  \item Choose an EM subbundle built from a 2D transverse polarization plane in the
        \emph{3D physical space} (tangent to the brane), consistent with observations.
  \item Show that excitations involving the fourth embedding component are either gapped,
        confined to the electron sector, or otherwise dynamically suppressed.
\end{itemize}

% ============================================================
\section{Simulation Protocol (Aligned with Current Code)}
\label{sec:simulation}

\subsection{Relevant code modules}
This checklist assumes the following modules and concepts exist as in the current project:
\begin{itemize}
  \item \texttt{branesim.core.state.BraneState}: positions/velocities/accelerations, rest positions.
  \item \texttt{branesim.core.grid.BraneGrid}: topology, neighbor indices, spatial coordinates.
  \item \texttt{branesim.core.solver.VelocityVerletSolver}: integration and energy diagnostics.
  \item \texttt{branesim.berry.analysis}: \texttt{compute\_berry\_time\_series} (time connection proxy).
  \item \texttt{branesim.electron.electron\_initialization}: rotating spherical-harmonic seed.
  \item Visualization utilities for displacement slices and videos.
\end{itemize}

\subsection{Electron initialization (rotating spherical-harmonic cavity seed)}
\label{subsec:electron_init}

This project initializes an ``electron'' as a localized, rotating standing-wave pattern on the 3D brane,
embedded in \(\R^4\). The initializer is designed to (i) create a controlled angular phase winding
(\(m\neq 0\)) in the intrinsic \((x,y,z)\) coordinates, (ii) realize that complex phase rotation in a
\emph{real-valued} simulator by using two quadratures, and (iii) ensure immediate coupling between the
spatial embedding components \((X^1,X^2,X^3)\) and the transverse embedding component \(X^4\) through
the 4D rest-length spring geometry.

\paragraph{Reference-space spherical coordinates.}
All mode functions are evaluated in the intrinsic/reference coordinates \(\x\in\Omega\subset\R^3\).
Given a chosen center \(\x_c\in\R^3\), define \(\mathbf{r}=\x-\x_c\), \(r=\|\mathbf{r}\|\), and
\[
\theta=\arccos\!\left(\frac{z}{r}\right),\qquad
\phi=\operatorname{atan2}(y,x),
\]
with the usual inverse mapping
\(
x=r\sin\theta\cos\phi,\;
y=r\sin\theta\sin\phi,\;
z=r\cos\theta.
\)

\paragraph{Angular factor \(Y_{\ell m}\).}
The angular structure is given by the spherical harmonic \(Y_{\ell m}(\theta,\phi)\), with standard
properties:
\begin{align}
\int_{S^2} Y_{\ell m}(\Omega)\,Y_{\ell' m'}^*(\Omega)\,d\Omega &= \delta_{\ell\ell'}\delta_{mm'},\\
\sum_{\ell=0}^{\infty}\sum_{m=-\ell}^{\ell} Y_{\ell m}(\Omega)\,Y_{\ell m}^*(\Omega')
&= \delta(\Omega-\Omega'),\\
Y_{\ell m}(\theta,\phi+2\pi) &= Y_{\ell m}(\theta,\phi),
\end{align}
and \(Y_{\ell m}\) is bounded and single-valued on \(S^2\) for integer \(\ell\ge 0\), \(|m|\le \ell\).
Choosing \(m\neq 0\) creates a nontrivial azimuthal phase winding \(\sim e^{im\phi}\), which is the
seed for a circulating phase gradient around the symmetry axis.

\paragraph{Radial factor \(R_{\ell n}\).}
We use a spherical-Bessel standing-wave profile
\[
R_{\ell n}(r)= j_\ell(k_{\ell n} r)\,w(r),
\qquad
k_{\ell n}=\frac{\alpha_{\ell n}}{a},
\]
where \(j_\ell\) is the spherical Bessel function, \(\alpha_{\ell n}\) is the \(n\)-th positive root of
\(j_\ell(\alpha)=0\), and \(a\) is the chosen electron radius (cavity radius) in intrinsic units.
The window \(w(r)\) is a smooth taper to avoid a hard discontinuity at \(r=a\); in code we use a
\(\tanh\)-ramp,
\[
w(r)=\frac12\left[1-\tanh\!\left(\frac{r-a}{a/s}\right)\right],
\]
with softness parameter \(s>0\) (larger \(s\) yields a softer edge).

\paragraph{Complex mode and real quadratures.}
Define the complex scalar mode
\[
\psi(\x)=R_{\ell n}(r)\,Y_{\ell m}(\theta,\phi),
\qquad
\psi(\x,t)=\psi(\x)\,e^{-i\omega t}.
\]
A real-valued integrator represents the rotating phase by embedding \(\Re\psi\) and \(\Im\psi\) into a
2D polarization plane in \(\R^4\). Let \(\mathbf{p}_1,\mathbf{p}_2\in\R^4\) be orthonormal vectors spanning this
plane (\(\mathbf{p}_i\cdot\mathbf{p}_j=\delta_{ij}\)). The initializer sets
\begin{align}
\uvec(\x,0) &= A\big(\Re\psi(\x)\,\mathbf{p}_1+\Im\psi(\x)\,\mathbf{p}_2\big),\\
\vvec(\x,0) &= A\,\omega\big(\Im\psi(\x)\,\mathbf{p}_1-\Re\psi(\x)\,\mathbf{p}_2\big),
\end{align}
which corresponds to a local circular rotation in the \((\mathbf{p}_1,\mathbf{p}_2)\) plane at angular frequency
\(\omega\). We choose
\[
\omega = c_{\mathrm{wave}}\,k_{\ell n},
\]
with \(c_{\mathrm{wave}}\) the brane wave speed in the (dimensionless) simulation calibration
(\(c_{\mathrm{wave}}=1\) in the current unit convention).

\paragraph{Amplitude normalization.}
The raw displacement field is scaled such that the maximum Euclidean norm of \(\uvec(\x,0)\) equals the
requested amplitude \(A\), i.e. \(A=\max_{\x}\|\uvec(\x,0)\|\) after normalization.

\paragraph{Polarization choices and \(X^4\) coupling.}
The code supports several fixed polarization presets (e.g. \texttt{xy}, \texttt{xz}, \texttt{yz},
\texttt{spatial}, \texttt{all}) and also accepts user-specified \(\mathbf{p}_1,\mathbf{p}_2\).
For electron seeding we typically use \texttt{spatial\_x4}, i.e. a plane that has support in the spatial
coordinates \emph{and} in \(X^4\). This ensures that \(X^4\) carries oscillatory content and momentum at
\(t=0\) and couples immediately to \((X^1,X^2,X^3)\) through the 4D rest-length spring geometry.

\paragraph{Optional containment scaffold in \(X^4\).}
In addition to the rotating seed, we optionally add a static Gaussian deformation in a chosen embedding
component (default: \(X^4\)),
\[
u_{\mathrm{well}}(r)=-D\exp\!\left(-\frac{r^2}{2\sigma^2}\right),
\]
added to \(X^{c_0}\) for some component index \(c_0\in\{1,2,3,4\}\).
This ``well'' is an initial geometric deformation (not an external potential). Any such deformation
changes neighbor distances in \(\R^4\) and therefore feeds back into the dynamics via the same rest-length
elasticity law that governs the full brane.


\paragraph{Initializer parameter set (implementation-facing).}
The electron seed is fully specified by the following parameters:
\begin{itemize}
  \item \(\ell,m,n\): spherical-harmonic degree/order and radial root index.
  \item \(a\): electron radius (intrinsic cavity radius) used in \(k_{\ell n}=\alpha_{\ell n}/a\).
  \item \(\x_c\): electron center in intrinsic coordinates.
  \item \(A\): target displacement amplitude after normalization.
  \item \(c_{\mathrm{wave}}\): wave speed used to set \(\omega=c_{\mathrm{wave}}k_{\ell n}\) (unity in current calibration).
  \item Polarization choice: either a preset plane (e.g. \texttt{spatial\_x4}) or explicit orthonormal vectors
        \(\mathbf{p}_1,\mathbf{p}_2\in\R^4\).
  \item \(s\): taper softness used in the window \(w(r)\) at the edge \(r\approx a\).
  \item Optional containment scaffold parameters: component index \(c_0\), depth \(D\), and width \(\sigma\) of the
        Gaussian well added to \(X^{c_0}\).
  \item Flags: whether to set initial velocities from the rotating quadrature formula (\texttt{set\_velocities}) and
        whether to rescale the seed to match the requested amplitude (\texttt{normalize}).
\end{itemize}

\paragraph{Design intent and current limitations.}
This initialization provides a reproducible, parameterized starting configuration with a controllable
azimuthal winding number \(m\) and a well-defined rotating phase in a real-valued simulation.
A long-lived confined ``electron'' would correspond to a self-consistent nonlinear eigenmode of the
coupled 4D rest-length elasticity system. The initializer alone does not enforce that equilibrium; it
provides the candidate initial condition from which such equilibria can be sought and tested.


\subsection{Checklist A: Data to record (minimal)}
\begin{enumerate}
  \item Record frames at stride \(s\): \(\X_p(t_k)\), \(\dot{\X}_p(t_k)\), and \(t_k\).
  \item Store \(\X_{0,p}\) (rest positions) to compute displacements \(\uvec_p = \X_p-\X_{0,p}\).
  \item Decide the embedding components \(E\) used for Berry analysis:
        \begin{itemize}
          \item EM: typically \(E=2\) (a transverse polarization plane in 3D space).
          \item Electron: \(E=2\) or \(E=4\) depending on whether the internal plane includes \(X^4\)
                (\texttt{polarization="spatial\_x4"} in the initializer).
        \end{itemize}
\end{enumerate}

\subsection{Checklist B: Temporal Berry phase / connection (already implemented)}
Use \texttt{branesim.berry.analysis.compute\_berry\_time\_series}:
\begin{enumerate}
  \item Build \(\uvec_k\in\R^{N\times E}\) and \(\vvec_k\in\R^{N\times E}\) per frame \(k\).
  \item Call \texttt{compute\_berry\_time\_series(frames\_u, frames\_v, times\_s)}.
  \item Save outputs: phase \(\gamma(t)\), increments \(\Delta\gamma\), connection proxy \(A_t\),
        amplitude and alpha mask for visualization.
\end{enumerate}

\subsection{Checklist C: Spatial Berry connection (to implement next)}
To validate Maxwell structure you need spatial components \(a_i(\x,t)\).
A robust discrete approach uses neighbor overlaps (gauge-invariant on loops):
\begin{enumerate}
  \item For each frame \(k\), compute normalized rays \(\psi_p\in\C^E\) at each lattice point \(p\).
  \item For each directed neighbor link \(p\to q\) (axis-aligned preferred):
        define discrete edge phase
        \[
        \theta_{p\to q} = \arg\!\big(\langle \psi_p \mid \psi_q\rangle\big).
        \]
        Then \(a_i \approx \theta/h\) for a step of length \(h\) along axis \(i\).
  \item Define a plaquette curvature (discrete \(f_{ij}\)) by summing edge phases around a loop:
        \[
        f_{ij} \approx \frac{1}{h^2}\Big(\theta_{p\to p+\hat{i}}+\theta_{p+\hat{i}\to p+\hat{i}+\hat{j}}
        -\theta_{p+\hat{j}\to p+\hat{i}+\hat{j}}-\theta_{p\to p+\hat{j}}\Big).
        \]
  \item Validate gauge invariance numerically: plaquette sum is invariant under
        \(\psi_p\to \e^{\ii\alpha_p}\psi_p\).
\end{enumerate}

\subsection{Checklist D: Maxwell-like validation tests}
Using \(f_{\mu\nu}\) from temporal and spatial differences:
\begin{enumerate}
  \item Define emergent fields: \(E_i \sim f_{i0}\), \(B_i \sim \tfrac12 \epsilon_{ijk} f_{jk}\).
  \item Check identities on the grid (finite differences):
        \(\nabla\cdot \mathbf{B}\approx 0\), and \(\partial_t \mathbf{B} + \nabla\times\mathbf{E}\approx 0\).
  \item For photon experiments, verify propagation speed and energy flux alignment.
  \item For electron experiments, verify localization and stability metrics:
        leakage, RMS radius, energy drift (matching existing diagnostics).
\end{enumerate}

\subsection{Checklist E: Mapping to existing experiment scripts}
\begin{itemize}
  \item Electron: \texttt{electron\_3d\_experiment.py} uses
        \texttt{initialize\_electron\_mode\_3d} with polarization options including \texttt{"spatial\_x4"}.
  \item Berry time series: use the recorded \(\uvec,\vvec\) frames from the experiment’s logging stride.
  \item Visualization: reuse displacement slice video pipeline; add Berry phase/alpha as an overlay.
\end{itemize}

% ============================================================
\section{Planned Results Section Structure}
\label{sec:results_plan}

\begin{enumerate}
  \item \textbf{Photon sector:} show two-polarization behavior and gauge-invariant curvature patterns.
  \item \textbf{Maxwell tests:} numerical residuals for Faraday and \(\nabla\cdot\mathbf{B}\).
  \item \textbf{Electron sector:} stability region over parameter sweeps (rest length, well depth, sigma, amplitude).
  \item \textbf{Internal holonomy:} demonstrate non-Abelian transport in a controlled deformation cycle (future work).
\end{enumerate}

% ============================================================
\section{Discussion and Failure Modes (state these early)}
\label{sec:discussion}

List what would falsify or strongly constrain the model:
\begin{itemize}
  \item Extra gapless polarizations that cannot be suppressed without fine tuning.
  \item Strong anisotropy / dispersion incompatible with experimental bounds.
  \item Inability to obtain a Maxwell-like effective action at leading order.
  \item Electron mode instability or inability to reproduce spin-like holonomy.
\end{itemize}

% ============================================================
\appendix
\section{Discrete Berry Curvature on a Lattice (derivation details)}
\label{app:discrete_curvature}

Provide a clean derivation of edge phases and plaquette sums, including handling of
undefined (low-amplitude) regions and numerical stabilization (overlap thresholds).

\section{Symbol Dictionary}
\label{app:dictionary}

\begin{tabular}{ll}
\(\Omega\) & intrinsic domain (periodic cell or bounded cube) \\
\(\X(\x,t)\) & embedding map \(\Omega\to\R^4\) \\
\(\X_0(\x)\) & rest configuration \\
\(\uvec=\X-\X_0\) & displacement field \\
\(\vvec=\partial_t \uvec\) & velocity field \\
\(\omega_0\) & carrier angular frequency (Compton-based in project) \\
\(\mathbf{a}\) & complex amplitude \(\sqrt{\omega_0}\uvec + \ii \vvec/\sqrt{\omega_0}\) \\
\(\ket{\psi}\) & normalized ray \(\mathbf{a}/\|\mathbf{a}\|\) \\
\(a_\mu\) & U(1) Berry connection \(\ii\braket{u}{\partial_\mu u}\) \\
\(f_{\mu\nu}\) & U(1) curvature \(\partial_\mu a_\nu-\partial_\nu a_\mu\) \\
\(\bbA_\mu\) & Wilczek--Zee connection (matrix) \\
\(\bbF_{\mu\nu}\) & Wilczek--Zee curvature \(\partial_\mu\bbA_\nu-\partial_\nu\bbA_\mu-\ii[\bbA_\mu,\bbA_\nu]\) \\
\(Y_{\ell m}\) & spherical harmonic on \(S^2\) \\
\(\ell,m\) & harmonic degree/order (\(|m|\le \ell\)) \\
\(j_\ell\) & spherical Bessel function (first kind) \\
\(\alpha_{\ell n}\) & \(n\)-th positive root of \(j_\ell(\alpha)=0\) \\
\(a\) & chosen electron radius (intrinsic cavity radius) \\
\(k_{\ell n}\) & radial wavenumber \(\alpha_{\ell n}/a\) \\
\(\omega\) & electron seed angular frequency \(c_{\mathrm{wave}}k_{\ell n}\) \\
\(\x_c\) & electron center in intrinsic coordinates \\
\(\mathbf{p}_1,\mathbf{p}_2\) & orthonormal polarization vectors in \(\R^4\) \\
\(A\) & displacement amplitude scale used for normalization \\
\(D,\sigma\) & Gaussian containment depth/width in embedding component \(X^{c_0}\) \\
\end{tabular}

% Uncomment if you use BibTeX:
% \bibliographystyle{unsrt}
% \bibliography{refs}

\end{document}
